\section{Launch-mass leverage}

The reason the closure cap matters economically is leverage. Because mass balance forces
$1-C$ kilograms of vitamins to be imported per kilogram the factory builds, the installed
mass per launched mass is exactly
\begin{equation}
L(C) \;=\; \frac{1}{1-C}. \label{eq:leverage}
\end{equation}
The function is unforgiving near one, as Fig.~\ref{fig:leverage} shows: at 67\% closure
each launched kilogram becomes about $3\times$ its mass in installed capacity, while at
97\% closure it becomes about $33\times$. The last few points of closure buy most of the
leverage, which is precisely why the electronics wall, by holding closure below one, caps
the payoff of the whole scheme.

\begin{figure}[!t]
  \centering
  \includegraphics[width=\columnwidth]{fig_leverage.pdf}
  \caption{Launch-mass leverage $1/(1-C)$ as a function of material closure $C$. Each
    kilogram launched becomes $1/(1-C)$ kilograms of installed factory, so closure of 0.67
    yields about $3\times$ while 0.97 yields about $33\times$, and the curve steepens
    sharply as $C$ approaches one, where the imported electronics hold it back.}
  \label{fig:leverage}
\end{figure}

The alternative to building in place is shipping a finished installation, and the rocket
equation makes that alternative exponentially expensive. The propellant mass ratio grows as
$\exp(\Delta v / v_e)$, so the cost of delivering finished mass to a destination rises
exponentially with the mission's delta-v \cite{tsiolkovsky-1903,sutton-biblarz-2016}. For
the representative budgets of an interplanetary mission, several kilometres per second of
delta-v against exhaust velocities of a few, the mass ratio is large
\cite{curtis-2020}, and even at modern launch prices the cost per delivered kilogram is
high \cite{jones-2018-launch-cost}. That exponential is the entire economic case for
self-replication: it is far cheaper to send a small, high-closure seed and let
(\ref{eq:leverage}) do the rest than to launch the installed capacity directly, the same
bootstrapping logic behind tonne-to-kilotonne space-industry studies
\cite{metzger-2013}. The electronics wall does not overturn this case; it bounds it. The
imported electronics that hold $C$ below one are exactly the mass that must still be flown,
at the rocket equation's exponential price, for every generation the factory builds.
